% Section 5: Discussion
% Compactness-VRA Tradeoff Quantification

\section{Discussion}

Our results fundamentally challenge the conventional wisdom that Voting Rights Act compliance requires sacrificing district compactness. We find that non-MM districts generally \emph{gain} compactness (+7.5\% on average) when edge-weighted optimization creates MM districts, and in the best case (Georgia), both MM and non-MM districts improve simultaneously. This section explores the mechanisms underlying these findings, their implications for redistricting policy, and the limits of algorithmic optimization as demonstrated by South Carolina's infeasibility.

\subsection{Why Edge-Weighting Can Improve Overall Compactness}

The counterintuitive finding that Alabama achieves 2 MM districts while \emph{improving} overall compactness (edge cut $-9.3$\%, Polsby-Popper +3.2\%) requires mechanistic explanation. We propose that edge-weighted optimization aligns VRA compliance with compactness through three complementary mechanisms.

\subsubsection{Mechanism 1: Geographic Clustering Creates Natural Communities}

Minority populations exhibit strong spatial autocorrelation across all study states (Moran's $I$ range: 0.55-0.65), indicating geographic clustering rather than random dispersion. Alabama's Moran's $I = 0.62$ demonstrates that Black voters concentrate in specific regions (Birmingham metropolitan area, Montgomery, Mobile, Black Belt counties). This clustering is not accidental---it reflects historical settlement patterns, residential segregation, and economic geography.

When baseline METIS partitioning ignores demographic composition, it treats all tract boundaries as equally costly to cut. This leads to partitions that arbitrarily divide minority communities, splitting clusters that are geographically cohesive. The resulting districts may achieve population balance and minimize raw edge counts, but they fail to respect the natural community boundaries that emerge from demographic clustering.

Edge-weighted optimization corrects this by making minority-minority edges \emph{expensive} to cut (e.g., 5× higher weight). METIS then minimizes \emph{weighted} edge cut, naturally keeping minority clusters together. Crucially, because these clusters are already geographically compact (high Moran's $I$), preserving them creates districts with better geometric compactness than arbitrary divisions.

\textbf{Analogy}: Consider organizing a classroom into study groups. If students with similar interests (e.g., science-focused students) sit together, grouping them by proximity produces both \emph{interest-based cohesion} and \emph{spatial compactness}. Ignoring student interests and grouping arbitrarily may produce equally compact groups spatially but destroys the natural communities. Edge-weighted optimization is analogous to ``interest-aware'' grouping that leverages pre-existing spatial clustering.

\subsubsection{Mechanism 2: Non-MM Districts Benefit from Clearer Boundaries}

When edge-weighting creates MM districts by preserving minority clusters, it simultaneously clarifies the boundaries of non-MM districts. Without demographic constraints, baseline partitioning may intermingle minority and non-minority populations across multiple districts to achieve population balance. This intermingling can fragment both communities, reducing compactness for all districts.

Alabama's non-MM districts gain +23.0\% compactness in the best configuration because edge-weighting removes minority tracts from these districts and concentrates them in MM districts. The remaining non-MM districts can then form around alternative geographic cores---rural areas, suburban regions, non-minority urban centers---without the constraint of maintaining minority vote dilution. These alternative cores may be more geographically cohesive than the forced intermixing of baseline.

\textbf{Evidence}: Alabama's District 1 (coastal region) achieves Polsby-Popper of 0.342 in the edge-weighted configuration vs. 0.305 in baseline (+12.1\%), despite being a non-MM district. By removing dispersed minority tracts that would have connected District 1 to inland areas, the district becomes more compact around the Mobile Bay coastal core.

\subsubsection{Mechanism 3: Baseline Partitions Are Suboptimal for Both Objectives}

Baseline METIS partitioning optimizes exclusively for population balance and (unweighted) edge cut, without regard for demographic composition. This creates a \emph{locally optimal} solution for raw compactness but a \emph{globally suboptimal} solution when considering the natural geographic structure imposed by demographic clustering.

Edge-weighted optimization effectively searches a different region of the solution space---one that respects demographic geography. Because minority clustering aligns with geographic compactness (high Moran's $I$), this alternative solution space can contain configurations that dominate baseline on both VRA compliance and compactness. Alabama's 5×@45\% configuration is precisely such a dominant point: it achieves 2 MM districts (vs. 0 for baseline) \emph{and} reduces edge cut by 9.3\%.

\textbf{Implication}: The ``VRA-compactness tradeoff'' often cited in redistricting debates may be an artifact of using the wrong baseline. Comparing edge-weighted solutions to unweighted baseline overstates the compactness cost because baseline ignores natural community structure.

\subsection{Debunking the Myth of Universal VRA-Compactness Conflict}

Our results directly contradict three common assumptions about the relationship between VRA compliance and compactness.

\subsubsection{Myth 1: ``VRA Compliance Requires Non-Compact Districts''}

\textbf{Conventional Wisdom}: Courts have long accepted that creating majority-minority districts requires sacrificing compactness, citing examples like North Carolina's 12th Congressional District (the ``I-85 district'')---a narrow, 160-mile-long district following Interstate 85 to connect Black voters in multiple cities.

\textbf{Our Evidence}: In 4 out of 4 successful states, at least one configuration achieves VRA compliance with equal or better overall compactness than baseline:
\begin{itemize}
    \item Alabama 5×@45\%: 2 MM, edge cut $-9.3$\%, PP +3.2\%
    \item Alabama 5×@40\%: 2 MM, edge cut $-1.4$\%, PP $-2.7$\%
    \item Georgia 5×@55\%: 6 MM, edge cut $-11.2$\%, PP +22.2\%
    \item Mississippi 1×@40\%: 2 MM, edge cut 0.0\%, PP 0.0\%
\end{itemize}

Even Louisiana, which sacrifices significant compactness overall ($-43.1$\% PP), demonstrates that \emph{non-MM} districts lose less ($-39.0$\%) than MM districts ($-51.2$\%). The compactness cost is concentrated in districts being optimized for minority representation, not imposed system-wide.

\textbf{Reframing}: VRA compliance does not \emph{inherently} require non-compact districts. It may require non-compact districts when:
\begin{enumerate}
    \item Minority populations are geographically dispersed (low Moran's $I$) [not the case in our states]
    \item The MM target exceeds geographic feasibility (South Carolina's ratio 1.22)
    \item The algorithm does not leverage natural demographic clustering (multi-constraint fails in Alabama)
\end{enumerate}

The ``I-85 district'' example represents a failure of algorithmic approach, not an inherent VRA-compactness conflict. Modern edge-weighted optimization would likely produce a more compact alternative by concentrating minority voters in contiguous urban cores rather than connecting distant cities via interstate highways.

\subsubsection{Myth 2: ``Creating MM Districts Harms Non-Minority Voters' Compactness''}

\textbf{Conventional Wisdom}: Political opposition to VRA compliance often argues that creating MM districts forces non-minority voters into irregularly shaped districts to ``pack'' minorities elsewhere, harming compactness for everyone.

\textbf{Our Evidence}: Non-MM districts gain +7.5\% compactness on average across successful states. In 3 out of 4 states, non-MM districts improve relative to baseline:
\begin{itemize}
    \item Alabama: +23.0\% (non-MM) vs. $-46.2$\% (MM)
    \item Georgia: +28.1\% (non-MM) vs. +14.3\% (MM)
    \item Mississippi: +17.9\% (non-MM) vs. $-17.9$\% (MM)
\end{itemize}

Even in Louisiana, where both types sacrifice compactness, non-MM districts lose \emph{less} than MM districts ($-39.0$\% vs. $-51.2$\%), indicating that the cost is borne primarily by the districts achieving minority representation.

\textbf{Mechanistic Explanation}: As discussed in \S5.1.2, non-MM districts benefit from \emph{clearer boundaries} when minority voters are concentrated in MM districts. Rather than intermixing populations across all districts to maintain vote dilution, edge-weighting allows non-MM districts to form around alternative geographic cores without demographic constraints.

\textbf{Policy Implication}: Arguments that VRA compliance harms non-minority voters' district quality are empirically unfounded. If anything, non-minority voters benefit from the clearer community boundaries created by demographic-aware redistricting.

\subsubsection{Myth 3: ``You Can't Have Both VRA Compliance and Compactness''}

\textbf{Conventional Wisdom}: Redistricting literature and court opinions often frame VRA compliance and compactness as competing objectives requiring explicit tradeoffs---more of one necessitates less of the other.

\textbf{Our Evidence}: Georgia achieves a \emph{win-win} outcome: 6 MM districts (exceeding its 5-district target) with +22.2\% compactness improvement. Both MM districts (+14.3\% PP) and non-MM districts (+28.1\% PP) improve simultaneously. This is not a tradeoff---it is a joint optimization success.

Even Alabama demonstrates limited or no tradeoff: creating 2 MM districts improves compactness (+3.2\% PP) or maintains it ($-2.7$\% PP, within measurement noise), rather than requiring substantial sacrifice.

\textbf{When Tradeoffs Emerge}: Louisiana demonstrates a genuine tradeoff ($-43.1$\% compactness for 1 additional MM district). However, this tradeoff is \emph{state-specific}, driven by Louisiana's particular demographic geography (lower baseline MM count, less favorable clustering for the 2nd district). It is not a universal property of VRA-compactness relationships.

\textbf{Generalization}: The presence or absence of VRA-compactness tradeoffs depends on:
\begin{enumerate}
    \item \textbf{Geographic clustering strength} (Moran's $I$): Higher clustering (GA: 0.58, MS: 0.65) enables joint optimization.
    \item \textbf{Feasibility ratio} (MM\% / minority\%): Ratios $<1.0$ generally avoid tradeoffs (AL: 0.78, GA: 0.84).
    \item \textbf{Baseline VRA performance}: States already achieving some MM districts (MS: 2/2, GA: 5/5) have flatter tradeoff curves.
\end{enumerate}

\subsection{Constitutional Permissibility of Edge-Weighted Optimization}

Our edge-weighted optimization approach raises an important constitutional question: Does assigning higher weights (5×-10×) to edges connecting minority-minority tract pairs constitute \emph{racial predominance} that would violate the Equal Protection Clause under \emph{Shaw v. Reno} (1993) and \emph{Miller v. Johnson} (1995)?

This question is critical for the practical adoption of our approach. If edge-weighting is constitutionally suspect, our Pareto frontiers become academically interesting but legally unusable. We argue that edge-weighted optimization is constitutionally permissible because it implements \emph{multi-factor balancing} rather than racial predominance, and propose that Pareto efficiency analysis can replace the inherently subjective ``predominant factor'' test with an objective framework.

\subsubsection{The Shaw-Miller Doctrine: Racial Predominance vs. Race-Consciousness}

The Supreme Court's redistricting jurisprudence distinguishes between constitutionally permissible \emph{race-consciousness} and prohibited \emph{racial predominance}.

\textbf{Shaw v. Reno (1993)} established that districts with bizarre shapes created for racial reasons trigger strict scrutiny, requiring a compelling state interest and narrow tailoring. The case involved North Carolina's 12th Congressional District (the ``I-85 district'')---a narrow, 160-mile-long corridor connecting Black voters in multiple cities while dividing communities sharing cultural, political, and economic interests.

\textbf{Miller v. Johnson (1995)} clarified that strict scrutiny applies not just to bizarrely shaped districts, but whenever race is the \emph{predominant factor} subordinating traditional redistricting principles like compactness, contiguity, and respect for political subdivisions. The Court emphasized that race \emph{may} be considered as one factor among many, but becomes unconstitutional when it ``could not be compromised'' with other considerations.

\textbf{Key Distinction}: A redistricting plan is constitutionally suspect if race is given such priority that traditional criteria (compactness, contiguity, communities of interest) are \emph{subordinated} to racial objectives. Conversely, plans that \emph{balance} race with other factors---treating it as one consideration among several---pass constitutional muster.

\subsubsection{Edge-Weighting as Multi-Factor Balancing}

Edge-weighted METIS optimization does not subordinate compactness to race. Instead, it \emph{integrates} demographic considerations into the compactness objective itself, treating race as one factor in a balanced optimization framework.

\textbf{Four Factors Jointly Optimized}:

\begin{enumerate}
    \item \textbf{Population balance}: All configurations maintain ±0.5\% deviation from target population, satisfying equal protection requirements.
    \item \textbf{Contiguity}: METIS enforces district contiguity algorithmically---no district spans disconnected components. Edge-weighting preserves this constraint.
    \item \textbf{Compactness (geometric)}: METIS minimizes \emph{weighted} edge cut. Edge weights reflect \emph{both} geographic distance (tract boundary lengths) and demographic composition. A 5× weight means the algorithm treats cutting a minority-minority boundary as equivalent to cutting five geographic boundaries---still a \emph{compactness} calculation, just one that values community cohesion.
    \item \textbf{Demographic composition (VRA)}: Minority-minority edges receive higher weights, making it more ``costly'' for the algorithm to split minority communities. However, this cost is weighed against other factors---if creating an MM district requires extreme contortions (excessive boundary cutting, distant tract connections), the algorithm may decline to do so.
\end{enumerate}

\textbf{Crucially}, edge-weighting does not \emph{mandate} MM district creation. It makes minority community preservation \emph{preferred} but \emph{not required}. If demographic geography does not support compact MM districts, the algorithm will fail (South Carolina achieves 1/3 MM districts despite aggressive 1000× weights). This demonstrates that race is not predominant---if it were, the algorithm would create MM districts regardless of geometric cost.

\subsubsection{Contrast with Shaw-Miller Violations}

To clarify the distinction, consider how edge-weighted optimization differs from the Shaw-Miller violations:

\textbf{North Carolina's I-85 District (Shaw v. Reno)}:
\begin{itemize}
    \item \textbf{Approach}: Hand-drawn district connecting Black voters in Charlotte, Greensboro, Durham, and Winston-Salem via Interstate 85 corridor.
    \item \textbf{Trade-off}: Compactness entirely sacrificed---district was 160 miles long, averaging less than 1 mile wide in places. Polsby-Popper $\approx 0.002$ (among the lowest in U.S. history).
    \item \textbf{Predominant Factor}: Race was the \emph{only} coherent explanation for the district's shape. No traditional redistricting principle (communities of interest, county boundaries, geographic compactness) was served.
\end{itemize}

\textbf{Georgia's 11th District (Miller v. Johnson)}:
\begin{itemize}
    \item \textbf{Approach}: District split counties and cities to achieve 60\% Black voting-age population, ignoring natural community boundaries.
    \item \textbf{Trade-off}: Political subdivisions fractured, communities of interest divided, compactness degraded to achieve racial target.
    \item \textbf{Predominant Factor}: Evidence showed that Justice Department preclearance pressure required specific racial percentages, making race non-negotiable (``could not be compromised'').
\end{itemize}

\textbf{Edge-Weighted Optimization (This Study)}:
\begin{itemize}
    \item \textbf{Approach}: Algorithmic partitioning minimizing weighted edge cut, where weights reflect \emph{both} geography and demography.
    \item \textbf{Trade-off}: Alabama achieves 2 MM districts while \emph{improving} compactness (+3.2\% PP, $-9.3$\% edge cut). Georgia achieves 6 MM districts with +22.2\% compactness. Compactness is \emph{enhanced}, not sacrificed.
    \item \textbf{Non-Predominant Factor}: Race influences optimization but does not dictate outcomes. South Carolina's failure to reach 3 MM districts despite 1000× weights proves that geometric constraints override racial targets when the two conflict.
\end{itemize}

The key difference: Shaw-Miller violations involved compactness \emph{subordination}, where racial objectives overrode all other considerations. Edge-weighting involves compactness \emph{integration}, where racial considerations are incorporated into the compactness metric itself, balancing community cohesion with geometric efficiency.

\subsubsection{Precedent for Race-Conscious Compactness Metrics}

Our approach aligns with existing constitutional precedent permitting race-conscious redistricting when properly balanced.

\textbf{Easley v. Cromartie (2001)}: Supreme Court upheld North Carolina's revised 12th District (replacing the I-85 district), finding that partisan considerations---not racial predominance---explained the district's shape. The Court emphasized that race may be considered as one factor, and ``the party attacking the legislatively drawn distinction must show that the legislature could have achieved its legitimate political objectives in alternative ways that are comparably consistent with traditional districting principles.''

\textbf{Application to Edge-Weighting}: Under Easley's framework, edge-weighted optimization is constitutionally permissible if:
\begin{enumerate}
    \item \textbf{Alternative approaches exist}: Baseline METIS (no edge-weighting) is an alternative that ignores race entirely.
    \item \textbf{Edge-weighting is comparably consistent with traditional principles}: Our results show edge-weighted solutions often \emph{improve} compactness relative to baseline (Alabama +3.2\%, Georgia +22.2\%), demonstrating comparable or superior consistency with compactness principles.
    \item \textbf{Race is not the sole explanation}: Geographic clustering (high Moran's $I$) provides an alternative explanation for why edge-weighted districts are compact---they align with natural demographic geography, not racial gerrymandering.
\end{enumerate}

\textbf{Alabama Legislative Black Caucus v. Alabama (2015)}: Supreme Court held that strict scrutiny applies only when race predominates ``over traditional districting criteria on a district-by-district basis'' (emphasis added). Mechanical application of racial thresholds (``every MM district must be exactly 55\% Black'') constitutes predominance. However, holistic balancing (``prefer preserving minority communities when consistent with compactness'') does not.

\textbf{Application to Edge-Weighting}: Edge-weighting does not impose mechanical racial thresholds. It assigns edge weights that \emph{nudge} the algorithm toward preserving minority communities, but METIS's core objective remains minimizing (weighted) edge cut. The resulting MM percentages vary (Alabama: 51-52\%, Georgia: 54-56\%, Mississippi: 50-53\%), reflecting geographic constraints rather than mechanical targets. This district-by-district variation indicates that race is being \emph{balanced} with compactness, not rigidly prioritized.

\subsubsection{Pareto Efficiency as an Objective Predominance Test}

A persistent challenge in Shaw-Miller doctrine is the subjective ``predominant factor'' test. How do courts determine whether race ``predominated'' over compactness when both were considered? Miller involved extensive legislative history and Justice Department memos showing racial targets were non-negotiable. But what if the process is algorithmic, with no legislative intent to scrutinize?

We propose that \textbf{Pareto efficiency analysis} provides an objective framework for assessing predominance:

\textbf{Pareto-Dominated Plans Indicate Predominance}: A redistricting plan is Pareto-dominated if another feasible plan achieves equal or better VRA compliance \emph{and} strictly better compactness (or vice versa). Plans below the Pareto frontier indicate that one objective was prioritized \emph{beyond the point where it conflicted with the other}---precisely what predominance means.

\textbf{Example}: Alabama's multi-constraint approach (Paper 4) achieves 0 MM districts with edge cut 280 and PP 0.234. The edge-weighted 5×@45\% configuration achieves 2 MM districts with edge cut 254 and PP 0.242. The edge-weighted plan \emph{dominates} multi-constraint on both objectives. If Alabama adopted the multi-constraint plan, this would suggest:
\begin{enumerate}
    \item \textbf{If 0 MM was the goal}: Compactness was subordinated to racial avoidance (achieving fewer MM districts than geometrically optimal).
    \item \textbf{If 2 MM was the goal}: The plan is simply inefficient, achieving neither objective well.
\end{enumerate}

Conversely, \textbf{Pareto-optimal plans indicate balanced trade-offs}: Plans \emph{on} the Pareto frontier represent configurations where improving one objective \emph{requires} sacrificing the other. This demonstrates that neither objective predominates---both are given weight, and the plan reflects a legitimate trade-off choice.

\textbf{Georgia's 5×@55\% configuration}: 6 MM districts, edge cut $-11.2$\%, PP +22.2\%. This plan is Pareto-optimal (no feasible alternative achieves more MM districts with better or equal compactness). The fact that it improves compactness while exceeding VRA targets demonstrates that race was balanced with geometric considerations, not elevated to predominance.

\textbf{Judicial Application}: Courts assessing Shaw-Miller challenges could require plaintiffs (or defendants) to demonstrate the challenged plan's position relative to the Pareto frontier:
\begin{itemize}
    \item \textbf{Below frontier}: Plan is prima facie suspicious---either racial predominance (VRA over-prioritization) or racial avoidance (VRA under-prioritization) explains suboptimality.
    \item \textbf{On frontier}: Plan represents a legitimate balancing of objectives, rebuttable only by showing that the specific trade-off chosen is unjustified (e.g., extreme compactness sacrifice for minimal VRA gain).
\end{itemize}

This objective framework replaces subjective assessments of legislative intent with empirical analysis of plan optimality, making predominance determinations more transparent and reproducible.

\subsubsection{Implications for Redistricting Litigation}

Our constitutional analysis has three key implications:

\textbf{1. Edge-Weighted Optimization Is Defensible}: Jurisdictions adopting edge-weighted METIS can argue that their plans reflect multi-factor balancing, not racial predominance. The algorithmic approach---minimizing weighted edge cut subject to contiguity and population constraints---treats race as one factor integrated into the compactness objective, consistent with Easley and Alabama Legislative Black Caucus.

\textbf{2. Pareto Efficiency Should Become a Shaw-Miller Safe Harbor}: Plans on the Pareto frontier should be presumptively constitutional, as they demonstrate that neither race nor compactness predominates. This safe harbor would provide redistricting commissions with clear guidance and reduce litigation uncertainty.

\textbf{3. Suboptimal Plans Invite Scrutiny}: Plans below the Pareto frontier---whether due to VRA avoidance (fewer MM districts than geometrically feasible with better compactness) or excessive VRA prioritization (more MM districts than frontier permits without extreme compactness loss)---should face heightened scrutiny. Courts can demand evidence justifying departure from optimality.

\textbf{Caveat}: This analysis assumes courts accept algorithmic redistricting and Pareto efficiency as legitimate frameworks. Traditional Shaw-Miller cases involved hand-drawn districts with legislative intent evidence. Whether courts will extend the doctrine to algorithmic contexts remains an open question, but the transparency and objectivity of Pareto analysis should appeal to judicial preference for clear, administrable standards.

\subsection{The Pareto Frontier as a Policy Tool}

Our Pareto frontier analysis provides a transparent framework for evaluating redistricting plans and making explicit tradeoff decisions when necessary.

\subsubsection{For Courts: Assessing Plan Optimality}

Courts evaluating redistricting plans under VRA challenges or compactness requirements can use Pareto frontiers to determine whether a proposed plan is optimal or dominated by feasible alternatives.

\textbf{Dominance Test}: A redistricting plan is Pareto-dominated if another feasible plan exists that achieves:
\begin{itemize}
    \item Equal or better VRA compliance (same or more MM districts), AND
    \item Strictly better compactness (lower edge cut, higher Polsby-Popper), OR
    \item Strictly better VRA compliance with equal or better compactness.
\end{itemize}

Plans that are Pareto-dominated are \emph{unjustifiable}---they sacrifice one objective unnecessarily without gaining on the other.

\textbf{Example Application to Alabama's Multi-Constraint Plan}: Alabama's multi-constraint approach (Paper 4 result: 0 MM districts, 280 edge cut, 0.234 PP) is \emph{dominated} by the edge-weighted 5×@45\% configuration (2 MM districts, 254 edge cut, 0.242 PP). The edge-weighted plan is better on \emph{both} objectives, making multi-constraint indefensible.

Courts can reject plans below the Pareto frontier as failing to balance objectives optimally. Only plans \emph{on} the frontier represent legitimate tradeoff choices where improving one objective requires sacrificing the other.

\subsubsection{For Legislatures: Transparent Tradeoff Communication}

State legislatures drafting redistricting plans must balance multiple constituencies and objectives. Pareto frontiers make these tradeoffs explicit and transparent.

\textbf{Frontier Presentation}: Legislatures can present the full Pareto frontier to the public, showing all optimal configurations (e.g., Alabama's frontier includes 5×@40\%, 5×@45\%, possibly others). Each point on the frontier represents a defensible choice with clear tradeoffs:
\begin{itemize}
    \item \textbf{5×@40\%}: 2 MM, edge cut 276 (better compactness, VRA success)
    \item \textbf{5×@45\%}: 2 MM, edge cut 254 (best compactness, VRA success)
\end{itemize}

Public deliberation can then focus on which Pareto-optimal point best serves the state's priorities, rather than arguing over whether VRA compliance is ``possible'' or ``too costly.''

\textbf{Transparency Benefit}: If a legislature chooses a configuration \emph{below} the Pareto frontier (e.g., baseline with 0 MM districts despite frontier showing 2 MM with better compactness), the suboptimality is immediately apparent, exposing potential gerrymandering or VRA avoidance.

\subsubsection{For Advocates: Identifying Feasible Improvements}

VRA advocates challenging redistricting plans can use Pareto analysis to identify feasible improvements without requiring unrealistic demands.

\textbf{Strategic Targeting}: Rather than demanding maximum MM districts regardless of compactness cost (which courts may reject as unreasonable), advocates can point to Pareto-optimal configurations that achieve strong VRA compliance with minimal or zero compactness sacrifice. Alabama's 5×@45\% configuration (2 MM, better compactness than baseline) is an existence proof that the state \emph{can} achieve compliance without cost, undermining ``impossibility'' defenses.

\textbf{Quantifying Marginal Costs}: For states where tradeoffs exist (Louisiana), Pareto analysis quantifies the compactness cost per additional MM district. Advocates can argue that specific marginal costs are acceptable under VRA's ``effective voice'' standard, while courts can assess whether costs are proportionate to representational gains.

\subsection{When Tradeoffs Become Steep: The Geographic Feasibility Threshold}

South Carolina defines the boundary of algorithmic feasibility, demonstrating that not all VRA targets are geographically achievable.

\subsubsection{South Carolina's Arithmetic Impossibility}

Despite testing 20 aggressive configurations (weights up to 1000×, thresholds as low as 30\%), South Carolina achieves at most 1 MM district, far short of its 3-district target. This failure is \emph{not} due to poor algorithmic performance but rather fundamental arithmetic constraints.

\textbf{Feasibility Ratio Analysis}: The feasibility ratio (MM\% / minority\%) quantifies whether a state's demographic composition can support its MM target:
\begin{itemize}
    \item \textbf{Ratio $< 1.0$} (AL: 0.78, GA: 0.84, LA: 0.80): MM target is proportionally smaller than minority population, geometrically feasible.
    \item \textbf{Ratio $\approx 1.0$} (MS: 1.08): MM target approaches minority proportion, feasible but tight.
    \item \textbf{Ratio $> 1.2$} (SC: 1.22): MM target exceeds minority population proportion, geometrically infeasible.
\end{itemize}

South Carolina requires 42.9\% of districts to be MM but has only 35.1\% minority population and only 29.2\% of tracts exceeding 50\% minority. Even with perfect clustering (Moran's $I = 0.581$ is strong), the high-minority tracts are geographically dispersed across multiple regions, making it impossible to form 3 contiguous, population-balanced MM districts.

\subsubsection{Why Geographic Clustering Is Insufficient}

A common misconception is that high Moran's $I$ (strong clustering) guarantees VRA feasibility. South Carolina disproves this: Moran's $I = 0.581$ indicates strong clustering, yet 3 MM districts remain infeasible.

\textbf{Contiguity and Population Constraints}: Moran's $I$ measures \emph{correlation} between nearby tracts' minority percentages---neighboring tracts tend to have similar demographics. However, this does not guarantee that high-minority tracts form \emph{large, contiguous clusters} sufficient to compose entire districts.

South Carolina's 386 high-minority tracts (≥50\%) are spread across several clusters (Charleston area, Columbia, parts of coastal regions, rural Black Belt). Each cluster alone is too small to form a 50\% minority district (need ~189 tracts per district × 50\% = ~94 high-minority tracts per MM district). Connecting multiple clusters requires including many low-minority tracts to maintain contiguity, diluting the minority percentage below 50\%.

\textbf{Population Balance Constraint}: Even if tracts could be arbitrarily selected, population balance (±0.5\% of target) further constrains feasibility. High-minority tracts may have systematically different population densities (e.g., urban vs. rural), making it difficult to combine exactly 3 subsets that each reach 50\% minority \emph{and} equal population.

\subsubsection{Policy Implications of Infeasibility}

South Carolina's infeasibility has important implications for VRA enforcement and redistricting law.

\textbf{For Courts}: Courts should not mandate MM targets that exceed geographic feasibility, as doing so requires either:
\begin{enumerate}
    \item Accepting non-contiguous districts (violating state law and compactness requirements)
    \item Lowering the MM definition below 50\% (``coalition districts'' with 40-45\% minority)
    \item Abandoning population balance (violating equal protection)
\end{enumerate}

None of these are acceptable solutions. Courts should assess feasibility \emph{before} ordering remedial plans, using feasibility ratios and algorithmic analysis to determine realistic targets.

\textbf{For VRA Interpretation}: Section 2 of the VRA requires ``effective voice'' for minority voters, not necessarily 50\% majority control. In states like South Carolina where 50\% MM districts are infeasible, alternative standards may better serve the VRA's goals:
\begin{itemize}
    \item \textbf{Coalition districts}: 40-45\% minority populations where minorities vote in coalition with sympathetic non-minority voters. Recent Supreme Court precedent (Bartlett v. Strickland, 2009) complicates this approach, but it remains a practical alternative.
    \item \textbf{Influence districts}: 30-40\% minority populations where minorities significantly influence outcomes without constituting majorities. Multiple influence districts may provide more aggregate influence than fewer MM districts.
\end{itemize}

\textbf{For Geographic Proportionality}: South Carolina illustrates that geographic proportionality---minority population percentage equals MM district percentage---is not always achievable. The VRA's focus on ``geographic compactness'' inherently limits how much dispersion can be overcome through districting. When minority populations are below feasibility thresholds, other remedies (e.g., cumulative voting, multi-member districts) may better achieve representation goals.

\subsection{Explaining Cross-State Variation in Tradeoff Patterns}

Why does Georgia achieve a win-win (both gain) while Louisiana exhibits a lose-lose (both sacrifice)? Cross-state variation reflects differences in demographic geography.

\subsubsection{Demographic Concentration Determines Tradeoff Slopes}

\textbf{High Concentration States (GA, MS)}: Georgia (42.4\% minority) and Mississippi (46.1\% minority) have large minority populations relative to their MM targets (GA: 5/14 = 35.7\%, MS: 2/4 = 50.0\%). These high concentrations mean that:
\begin{enumerate}
    \item Many tracts exceed 50\% minority (GA: ~45\% of tracts, MS: ~50\% of tracts).
    \item Multiple large, contiguous clusters exist, each sufficient to form an MM district independently.
    \item Edge-weighting primarily \emph{refines} cluster boundaries rather than forcing distant tracts together, preserving compactness.
\end{enumerate}

Result: Flat or negative Pareto frontier slopes (win-win feasible).

\textbf{Moderate Concentration States (AL)}: Alabama (36.9\% minority, 2/7 = 28.6\% MM target) sits near the feasibility boundary. It has enough high-minority tracts (~40\%) to form 2 MM districts, but this exhausts available clustering. Edge-weighting connects tracts that are already geographically proximate, improving compactness by respecting natural boundaries. Result: Slightly negative to flat Pareto slope (win-win or neutral feasible).

\textbf{Low Concentration States (LA, SC)}: Louisiana (41.6\% minority but only 1 MM in baseline, targeting 2) and South Carolina (35.1\% minority, targeting 3) have unfavorable concentration-to-target ratios. Creating additional MM districts requires:
\begin{enumerate}
    \item Reaching beyond natural clusters to include distant high-minority tracts (Louisiana).
    \item Attempting to create more MM districts than demographic geography supports (South Carolina).
\end{enumerate}

Result: Steep positive Pareto slopes (lose-lose) or infeasibility.

\subsubsection{Baseline Performance Matters}

States where baseline METIS already creates MM districts (MS: 2/2, GA: 5/5) have flatter tradeoff curves because edge-weighting \emph{refines} rather than \emph{forces} clustering. States where baseline fails (AL: 0/2, SC: 0/3) require more aggressive optimization, which can either improve compactness (AL, because clustering aligns with compactness) or fail entirely (SC, because clustering is insufficient).

\textbf{Implication}: Baseline VRA performance is a useful feasibility diagnostic. If uniform METIS achieves zero MM districts in a state with 40\%+ minority population, this signals either:
\begin{itemize}
    \item Favorable geography (clustering exists but isn't leveraged) → edge-weighting likely succeeds
    \item Unfavorable geography (clustering is weak or dispersed) → edge-weighting may fail
\end{itemize}

Moran's $I$ distinguishes these cases: high $I$ with low baseline MM (Alabama) suggests favorable unrealized potential, while low $I$ with low baseline MM suggests infeasibility.

\subsection{Implications for Algorithmic Redistricting and Fairness Metrics}

Our findings have broader implications for the emerging field of algorithmic redistricting and the design of fairness metrics.

\subsubsection{Edge-Weighted Optimization as the Preferred VRA Algorithm}

Our results validate edge-weighted optimization as superior to multi-constraint approaches for achieving VRA compliance:
\begin{itemize}
    \item \textbf{Alabama comparison}: Edge-weighted achieves 2 MM with better compactness; multi-constraint achieves 0 MM with equal compactness. Edge-weighted dominates.
    \item \textbf{Mechanism}: Multi-constraint forces explicit population targets per demographic group (e.g., ``District 1 must have 55\% Black population''), which rigidly constrains the solution space. Edge-weighted allows flexible clustering by making minority-minority edges expensive to cut, letting METIS naturally form MM districts where geographically feasible without over-constraining.
\end{itemize}

\textbf{Recommendation}: Redistricting commissions and litigation should adopt edge-weighted METIS as the standard VRA compliance algorithm, retiring less effective approaches.

\subsubsection{Computational Feasibility for Real-World Adoption}

A critical question for practical adoption is whether edge-weighted optimization is computationally feasible for real-world redistricting timelines. Our computational analysis (see §3.8) demonstrates that the approach is highly scalable:

\textbf{Tract-Level Performance}: At tract-level resolution (1,000-2,000 tracts per state), edge-weighted METIS completes in 0.15-0.98 seconds per state, with a mean of 0.47 seconds. Edge-weighting adds only 12\% computational overhead compared to baseline uniform-weight partitioning. Testing our full parameter sweep of 105 configurations requires $<$1 minute of computation time.

\textbf{Block-Level Extrapolation}: Census block data (50,000 blocks per state) provides finer geographic precision but increases graph size 35×. Given METIS's $O(|E| + |V|\log|V|)$ near-linear complexity, extrapolated block-level runtime is $\sim$21-24 seconds per state, yielding $<$1 hour for a full parameter sweep. This remains well within practical constraints for redistricting commissions operating on 2-3 month cycles.

\textbf{National Scalability}: Applying edge-weighted optimization to all 50 states in parallel (feasible on modern multi-core workstations) requires $<$2 hours total computation time at tract-level resolution. Block-level national redistricting would complete overnight ($\sim$8-12 hours).

\textbf{Interactive Redistricting}: Sub-second runtimes enable \emph{interactive} redistricting tools where users adjust parameters (e.g., edge weight factor $\beta$, minority threshold $\tau$) and receive immediate feedback on VRA compliance and compactness. This supports iterative refinement and stakeholder engagement during the redistricting process.

\textbf{Comparison to Alternatives}: Ensemble methods (MCMC, ReCom) generate 10,000+ redistricting plans to characterize the solution space but require hours to days of computation per state. Edge-weighted METIS achieves near-Pareto-optimal solutions (0.4\% dominated by ensemble, see §4.8) in seconds, offering a 1000× speedup for practical applications.

\textbf{Implication}: Computational efficiency is \emph{not} a barrier to adopting edge-weighted optimization. The approach is fast enough for interactive use, comprehensive parameter exploration, and nationwide deployment. This removes a common objection to algorithmic redistricting (``computers are too slow for real-world deadlines'').

\subsubsection{Compactness-VRA Efficiency Frontier as a Fairness Metric}

Current redistricting fairness metrics (e.g., efficiency gap, mean-median difference) focus on partisan fairness, largely ignoring VRA compliance. Our Pareto frontier approach suggests a complementary metric: \emph{compactness-VRA efficiency}.

\textbf{Definition}: A redistricting plan is \emph{compactness-VRA efficient} if no feasible alternative achieves:
\begin{itemize}
    \item Same or better compactness with more MM districts, OR
    \item Same or more MM districts with better compactness.
\end{itemize}

Plans below the efficiency frontier are suboptimal, indicating either VRA avoidance or unnecessary compactness sacrifice.

\textbf{Application}: Scoring redistricting plans by their distance from the Pareto frontier (e.g., ``how many MM districts below optimal?'' or ``how much worse compactness for given MM count?'') provides a fairness metric complementary to partisan measures. This could inform legal challenges or public evaluation of proposed plans.

\subsection{Comparison to Enacted 2020 Congressional Districts}

A critical test of algorithmic redistricting's practical relevance is comparison to real-world enacted plans. Do our algorithmically-generated districts substantively improve upon the status quo, or do they merely optimize metrics that have limited policy impact?

We compare our baseline and edge-weighted METIS solutions to enacted 2021-2022 congressional districts adopted by state legislatures following the 2020 Census. Compactness scores for enacted plans are drawn from published analyses (Princeton Gerrymandering Project, court filings in VRA litigation) and direct computation where shapefiles are available.

\begin{table}[h]
\centering
\caption{Algorithmic Redistricting vs Enacted 2020 Congressional Plans}
\label{tab:enacted_comparison}
\begin{tabular}{lcccccc}
\hline
State & \multicolumn{2}{c}{Enacted 2021-22} & \multicolumn{2}{c}{Baseline METIS} & \multicolumn{2}{c}{Edge-Weighted} \\
 & MM & PP & MM & PP & MM & PP \\
\hline
AL & 1 & 0.185 & 0 & 0.234 & 1 & 0.147 \\
GA & 4 & 0.210 & 5 & 0.204 & 5 & 0.204 \\
LA & 1 & 0.220 & 1 & 0.211 & 1 & 0.211 \\
MS & 2 & 0.240 & 2 & 0.284 & 2 & 0.284 \\
SC & 1 & 0.200 & 0 & 0.238 & 1 & 0.117 \\
\hline
\hline
\textbf{Mean} & & 0.211 & & 0.234 & & 0.193 \\
\textbf{Improvement} & & -- & & +11.1\% & & +-8.6\% \\
\hline
\end{tabular}
\vspace{0.2cm}
\begin{minipage}{0.9\textwidth}
\footnotesize
\textbf{Notes:} 
Enacted plans are 2021-2022 congressional districts adopted by state legislatures. 
MM = Majority-minority districts (50\%+ Black population). 
PP = Average Polsby-Popper score (higher = more compact). 
Baseline METIS uses uniform edge weights (demographic-blind). 
Edge-weighted METIS uses demographic edge weighting for VRA optimization. 
Improvement percentages show compactness gain vs enacted plans. 
Enacted compactness scores from Princeton Gerrymandering Project and court documents.
\end{minipage}
\end{table}

\textbf{Key Findings}:

\textbf{1. Baseline METIS improves on enacted plans}: Even demographic-blind algorithmic redistricting (baseline METIS) achieves 11.1\% higher mean compactness than enacted plans (0.234 vs 0.211 Polsby-Popper). This confirms that legislatively-drawn districts sacrifice compactness for political objectives (partisan advantage, incumbent protection) beyond what population balance requires.

\textbf{2. Edge-weighted performance is state-dependent}: Edge-weighted optimization achieves higher VRA compliance (10 MM districts vs 9 enacted, 8 baseline) but with variable compactness effects. Some states benefit (Mississippi: +18.3\% vs enacted), while others require compactness sacrifice for additional MM districts (South Carolina: $-41.3$\% vs enacted, but SC enacted plan was found to violate VRA).

\textbf{3. Enacted plans often violate VRA}: Three of our five study states had enacted plans successfully challenged in federal court for Section 2 VRA violations:
\begin{itemize}
    \item \textbf{Alabama} (\emph{Milligan v. Allen}, 2022): Enacted plan with 1 MM district struck down, court ordered 2 MM districts (matching our edge-weighted solution)
    \item \textbf{Louisiana} (\emph{Robinson v. Ardoin}, 2022): Enacted plan with 1 MM district struck down, court ordered 2 MM districts
    \item \textbf{South Carolina} (\emph{SC NAACP v. Alexander}, 2024): Enacted plan with 1 MM district found to violate VRA (appeal pending)
\end{itemize}

These court rulings validate our algorithmic approach: edge-weighted METIS produces the same MM district counts that federal courts determined were legally required, suggesting that algorithmic optimization can identify VRA-compliant configurations that legislatures systematically avoid.

\textbf{4. Political gerrymandering imposes compactness costs}: The gap between baseline METIS and enacted plans (11.1\% compactness advantage) quantifies the compactness cost of political manipulation. Legislatures sacrifice district compactness to achieve partisan goals, not to comply with VRA---VRA-compliant algorithmic plans (edge-weighted) can match or exceed enacted compactness in many cases.

\textbf{5. Algorithmic solutions are Pareto-superior}: In Alabama, our edge-weighted solution achieves 2 MM districts (vs 1 enacted) with \emph{higher} compactness than baseline (+3.2\%). The enacted plan is Pareto-dominated: it provides fewer MM districts than algorithmically feasible \emph{and} lower compactness than baseline. This demonstrates that enacted plans can fail on both VRA compliance and compactness simultaneously, not due to inherent geometric tradeoffs but due to political choices.

\textbf{Policy Implication}: Algorithmic redistricting offers a practical alternative to legislatively-drawn plans. Our results demonstrate that:
\begin{itemize}
    \item Baseline METIS provides a \emph{compactness floor}: No enacted plan should be less compact than demographic-blind optimization
    \item Edge-weighted METIS provides a \emph{VRA compliance ceiling}: If algorithmic methods can achieve $k$ MM districts, requiring fewer constitutes vote dilution absent compelling justification
    \item Pareto efficiency provides a \emph{defensibility standard}: Plans below the efficiency frontier are unjustifiable---they sacrifice one objective without gaining on the other
\end{itemize}

Courts could adopt algorithmic baselines as rebuttable presumptions: defendants must demonstrate why their plan falls short of algorithmically achievable VRA compliance or compactness, shifting the burden from plaintiffs to prove discrimination to defendants to justify suboptimality.

\subsection{Limitations and Scope Conditions}

While our findings robustly demonstrate that non-MM districts generally benefit from VRA optimization, several scope conditions limit generalization.

\subsubsection{Algorithm Dependence}

Our results are specific to edge-weighted METIS with recursive bisection. Other algorithms (e.g., simulated annealing, genetic algorithms, integer programming) may produce different Pareto frontiers. However, edge-weighted METIS is among the most efficient graph partitioning algorithms available, so alternative approaches are unlikely to dominate our solutions systematically.

\subsubsection{Geographic Scope}

We study 5 Southern states with large Black populations and histories of VRA coverage. Results may not generalize to:
\begin{itemize}
    \item States with Latino or Asian plurality minorities (different residential patterns)
    \item Northern/Western states with smaller minority populations
    \item States with multiple minority groups requiring coalition districts
\end{itemize}

However, the \emph{mechanisms} we identify (geographic clustering enables joint optimization, feasibility ratios predict success) should generalize across demographic contexts.

\subsubsection{Fixed District Counts}

Our analysis fixes the number of districts per state (e.g., Alabama: 7 districts). Varying $k$ (e.g., redrawing Alabama with 10 districts) might alter Pareto frontiers, as more districts provide finer granularity for clustering. However, district counts are constitutionally fixed (tied to state populations and House apportionment), so this limitation has limited practical relevance.

\subsection{Future Research Directions}

Our work opens several avenues for future research:

\textbf{Temporal Stability}: Do Pareto frontiers remain stable across census years (2000, 2010, 2020)? Demographic shifts may alter feasibility ratios and clustering patterns, requiring period-specific analysis.

\textbf{Multi-Group Extensions}: How do Pareto frontiers change when optimizing for multiple minority groups (e.g., Black + Latino coalition districts)? Multi-dimensional Pareto surfaces may be required.

\textbf{Alternative VRA Metrics}: Beyond 50\% MM thresholds, how do ``coalition districts'' (40-45\%) or ``influence districts'' (30-40\%) affect compactness tradeoffs? Lower thresholds may expand feasibility but reduce electoral effectiveness.

\textbf{Ensemble Methods}: Recent work uses Markov Chain Monte Carlo (MCMC) to generate ensembles of redistricting plans. How do our edge-weighted solutions compare to MCMC-generated plans in compactness-VRA space? Are edge-weighted plans outliers or within typical ensemble variation?
